\documentclass[11pt]{article}
\usepackage{color}
%\input{rgb}

%----------Packages----------
\usepackage{amsmath}
\usepackage{amssymb}
\usepackage{amsthm}
%\usepackage{amsrefs}
\usepackage{dsfont}
\usepackage{mathrsfs}
\usepackage{stmaryrd}
\usepackage[all]{xy}
\usepackage[mathcal]{eucal}
\usepackage{verbatim}  %%includes comment environment
\usepackage{fullpage}  %%smaller margins
%----------Commands----------

%%penalizes orphans
\clubpenalty=9999
\widowpenalty=9999





%% bold math capitals
\newcommand{\bA}{\mathbf{A}}
\newcommand{\bB}{\mathbf{B}}
\newcommand{\bC}{\mathbf{C}}
\newcommand{\bD}{\mathbf{D}}
\newcommand{\bE}{\mathbf{E}}
\newcommand{\bF}{\mathbf{F}}
\newcommand{\bG}{\mathbf{G}}
\newcommand{\bH}{\mathbf{H}}
\newcommand{\bI}{\mathbf{I}}
\newcommand{\bJ}{\mathbf{J}}
\newcommand{\bK}{\mathbf{K}}
\newcommand{\bL}{\mathbf{L}}
\newcommand{\bM}{\mathbf{M}}
\newcommand{\bN}{\mathbf{N}}
\newcommand{\bO}{\mathbf{O}}
\newcommand{\bP}{\mathbf{P}}
\newcommand{\bQ}{\mathbf{Q}}
\newcommand{\bR}{\mathbf{R}}
\newcommand{\bS}{\mathbf{S}}
\newcommand{\bT}{\mathbf{T}}
\newcommand{\bU}{\mathbf{U}}
\newcommand{\bV}{\mathbf{V}}
\newcommand{\bW}{\mathbf{W}}
\newcommand{\bX}{\mathbf{X}}
\newcommand{\bY}{\mathbf{Y}}
\newcommand{\bZ}{\mathbf{Z}}

%% blackboard bold math capitals
\newcommand{\bbA}{\mathbb{A}}
\newcommand{\bbB}{\mathbb{B}}
\newcommand{\bbC}{\mathbb{C}}
\newcommand{\bbD}{\mathbb{D}}
\newcommand{\bbE}{\mathbb{E}}
\newcommand{\bbF}{\mathbb{F}}
\newcommand{\bbG}{\mathbb{G}}
\newcommand{\bbH}{\mathbb{H}}
\newcommand{\bbI}{\mathbb{I}}
\newcommand{\bbJ}{\mathbb{J}}
\newcommand{\bbK}{\mathbb{K}}
\newcommand{\bbL}{\mathbb{L}}
\newcommand{\bbM}{\mathbb{M}}
\newcommand{\bbN}{\mathbb{N}}
\newcommand{\bbO}{\mathbb{O}}
\newcommand{\bbP}{\mathbb{P}}
\newcommand{\bbQ}{\mathbb{Q}}
\newcommand{\bbR}{\mathbb{R}}
\newcommand{\bbS}{\mathbb{S}}
\newcommand{\bbT}{\mathbb{T}}
\newcommand{\bbU}{\mathbb{U}}
\newcommand{\bbV}{\mathbb{V}}
\newcommand{\bbW}{\mathbb{W}}
\newcommand{\bbX}{\mathbb{X}}
\newcommand{\bbY}{\mathbb{Y}}
\newcommand{\bbZ}{\mathbb{Z}}

%% script math capitals
\newcommand{\sA}{\mathscr{A}}
\newcommand{\sB}{\mathscr{B}}
\newcommand{\sC}{\mathscr{C}}
\newcommand{\sD}{\mathscr{D}}
\newcommand{\sE}{\mathscr{E}}
\newcommand{\sF}{\mathscr{F}}
\newcommand{\sG}{\mathscr{G}}
\newcommand{\sH}{\mathscr{H}}
\newcommand{\sI}{\mathscr{I}}
\newcommand{\sJ}{\mathscr{J}}
\newcommand{\sK}{\mathscr{K}}
\newcommand{\sL}{\mathscr{L}}
\newcommand{\sM}{\mathscr{M}}
\newcommand{\sN}{\mathscr{N}}
\newcommand{\sO}{\mathscr{O}}
\newcommand{\sP}{\mathscr{P}}
\newcommand{\sQ}{\mathscr{Q}}
\newcommand{\sR}{\mathscr{R}}
\newcommand{\sS}{\mathscr{S}}
\newcommand{\sT}{\mathscr{T}}
\newcommand{\sU}{\mathscr{U}}
\newcommand{\sV}{\mathscr{V}}
\newcommand{\sW}{\mathscr{W}}
\newcommand{\sX}{\mathscr{X}}
\newcommand{\sY}{\mathscr{Y}}
\newcommand{\sZ}{\mathscr{Z}}

\newcommand{\id}{\mathsf{id}}
\newcommand{\tr}{\mathrm{tr}}
\newcommand{\Char}{\mathrm{char}}

\renewcommand{\phi}{\varphi}

%\renewcommand{\emptyset}{\O}

\providecommand{\abs}[1]{\lvert #1 \rvert}
\providecommand{\norm}[1]{\lVert #1 \rVert}


\providecommand{\x}{\times}




\providecommand{\ar}{\rightarrow}
\providecommand{\arr}{\longrightarrow}



\newcommand{\head}[1]{
	\begin{center}
		{\large #1}
		\vspace{.2 in}
	\end{center}
	
	\bigskip 
}



%----------Theorems----------

\newtheorem{theorem}{Theorem}[section]
\newtheorem{proposition}[theorem]{Proposition}
\newtheorem{lemma}[theorem]{Lemma}
\newtheorem{corollary}[theorem]{Corollary}

\theoremstyle{definition}
\newtheorem{definition}[theorem]{Definition}
\newtheorem*{definition*}{Definition}
\newtheorem{nondefinition}[theorem]{Non-Definition}
\newtheorem{exercise}[theorem]{Exercise}



\numberwithin{equation}{subsection}


%----------Title-------------
\newcommand{\hide}[1]{{\color{red} #1}} 
\newcommand{\com}[1]{{\color{blue} #1}} 
\newcommand{\meta}[1]{{\color{green} #1}} 

\begin{document}

\pagestyle{plain}


%%---  sheet number for theorem counter
\setcounter{section}{1}   

\head{Abstract Linear Algebra, Sections 41\&51, Spring 2025\\ HOMEWORK 7} 

This homework is due to {\bf December 3}, midnight. The total grade is 30 points. Please upload your solution on Canvas/Gradescope.

\begin{enumerate}

\item Let $A\in M_{n\times n}(\bbF)$ be a square matrix. The \emph{trace} $\tr(A)$ of $A$ is the sum
$$\tr(A) = \sum_{i=1}^n a_{ii} =a_{11}+ a_{22}+\ldots + a_{nn} $$
of its diagonal enrtries.
\begin{enumerate}
\item (\emph{1 point}) Find a $2\times 2$-matrix $A \in M_{2\times 2}(\bbF)$ such that $\tr(A^2)\neq \tr(A)^2$.
\item[] $A=\begin{bmatrix}
    0&1\\0&0
\end{bmatrix}\Rightarrow A^2=\begin{bmatrix}
    0&0\\0&0
\end{bmatrix}, \;\tr(A^2)=0\neq \tr(A)^2=1^2$.
\item (\emph{2 points}) Show that $\tr(AB) = \tr(BA)$ for any matrices $A \in M_{n\times m}(\bbF)$ and $B \in M_{m\times n}(\bbF)$. Note that $m$ is not necessary equal to $n$.
\item[] $\tr(AB)=\sum_{i=1}^n(AB)_{ii}=\sum_{i=1}^n\sum_{j=1}^mA_{ij}B{ji}\Rightarrow\text{Swap order of summation}$
\item[]$\Rightarrow\sum_{j=1}^m\sum_{i=1}^nB_{ji}A_{ij}=\sum_{j=1}^m(BA)_{jj}=\tr(BA)$. Thus $\tr(AB)=\tr(BA)$.
\item (\emph{2 points}) Find three matrices $A,B,C$ such that $\tr(ABC) \neq \tr(BAC)$.
\item[] $A=\begin{bmatrix}
    1&0\\0&0
\end{bmatrix},\; B=\begin{bmatrix}
    0&1\\0&0
\end{bmatrix},\; C=\begin{bmatrix}
    0&0\\1&0
\end{bmatrix}\Rightarrow ABC=\begin{bmatrix}
    1&0\\0&0
\end{bmatrix},\; BAC=\begin{bmatrix}
    0&0\\0&0
\end{bmatrix}.$
\item[] $\tr(ABC)=1\neq\tr(BAC)=0$.
\item (\emph{1 point}) Show that $\tr(C^{-1}AC)=\tr(A)$ for any square matrix $A \in M_{n\times n}(\bbF)$ and any invertible matrix $C \in M_{n\times n}(\bbF)$.
\item[] From (b): $\tr(C^{-1}AC)=\tr((AC)C^{-1})=\tr(A(CC^{-1}))=\tr(A\cdot I)=\tr(A)$.
\item (\emph{2 points}) Let $f\colon V \to V$ be an endomorphism and let $\mathcal{A}=\{v_1,\ldots, v_n\}$, $\mathcal{A}'=\{v'_1,\ldots, v'_n\}$ be bases of $V$. Show that
$$\tr([f]_{\mathcal{A}\mathcal{A}}) = \tr([f]_{\mathcal{A}'\mathcal{A}'}). $$
In other words, the trace of an endomorphism does not depend on the choice of a basis of $V$.
\item[] Suppose $P$ is the change of basis matrix from $\mathcal{A}$ to $\mathcal{A}'$. Then $\mathcal{A'}=P^{-1}\mathcal{A}P$. By part (d): $\tr(\mathcal{A}')=\tr(P^{-1}\mathcal{A}P)=\tr(\mathcal{A}).$ So trace is independent of basis choice.
\end{enumerate}

\item Let $A=\begin{bmatrix} 4 & 0 & 1 & 0\\ 2 & 2 & 3 & 0\\ -1 & 0 & 2 & 0\\ 4 & 0 & 1 &2\end{bmatrix}$.

\begin{enumerate}

\item (\emph{1 point}) Find all \emph{real} eigenvalues of $A$.
\item[] Fourth column is $(0,0,0,2)^T$. Thus, $A=\begin{bmatrix}
    M&0\\r^T&2
\end{bmatrix}$, where $M\in\bbR_{3\times 3}$, and $0$ is a $3\times 1$. Then, $\det(A-\lambda I)=\det(M-\lambda I)\cdot(2-\lambda)$. So one eigenvalue is 2. Compute $\det(M-\lambda I)$: $\begin{bmatrix}
    4-\lambda&0&1\\2&2-\lambda&3\\-1&0&2-\lambda
\end{bmatrix}\Rightarrow CP=(4-\lambda)((2-\lambda)^2)-1(-1(2-\lambda))=-\lambda^3+8\lambda^2-21\lambda+18\Rightarrow (2-\lambda)(\lambda-3)^2$. Hence eigenvalues over $\bbR$ are 2 (multiplicity 2) and 3 (multiplicity 2).
\item (\emph{2 points}) For each eigenvalue, $\lambda$, that you found in part (a), find a basis for the corresponding eigenspace, $E_{\lambda}$. 
\item[] For $\lambda=2$, solve $(A-2I)=\begin{bmatrix}
    2&0&1&0\\2&0&3&0\\-1&0&0&0\\4&0&1&0
\end{bmatrix}\Rightarrow\begin{cases}
    2x_1+1x_3=0\\2x_1+3x_3=0\\-1x_1=0\\4x_1+1x_3=0
\end{cases}$. $x_1=0\Rightarrow x_3=0$ Eigenspace, $E_2=\{(0,x_2,0,x_4): x_2,x_4\in\bbR\}=\text{span}\{(0,1,0,0)^T,(0,0,0,1)^T\}$.
\item[] For $\lambda=3$, solve $(A-3I)=\begin{bmatrix}
    1&0&1&0\\2&-1&3&0\\-1&0&-1&0\\4&0&1&-1
\end{bmatrix}\Rightarrow\begin{cases}
    1x_1+1x_3=0\\2x_1-x_2+3x_3=0\\-1x_1-1x_3=0\\4x_1+1x_3-1x_4=0
\end{cases}$, $x_3=-x_1$ Substitute into (2) $\Rightarrow 2x_1-x_2+3(-x_1)=0\Rightarrow -x_1-x_2=0\Rightarrow x_2=-x_1$
\item[] Substitute into (4) $\Rightarrow 4x_1+1(-x_1)-1x_4=0\Rightarrow 3x_1-1x_4=0\Rightarrow x_4=3x_1$
\item[] Eigenspace, $E_3=\{(x_1,-x_1,-x_1,3x_1):x_1\in\bbR\}=\text{span}\{(1,-1,-1,3)^T\}$.

\item (\emph{2 points}) If possible, diagonalize $A$ over $\mathbb{R}$. If this is not possible, then explain why not.
\item[] To be diagonalizable, the sum of dimensions of eigenspaces must equal the size of $A$ (4). Since $\dim E_2=2$ and $\dim E_3=1$, for a total of 3 linearly independent eigenvectors in the combined eigenspace, $A$ is not diagonalizable over $\bbR$. 

\end{enumerate}

\item Let $A=\begin{bmatrix} -1 & 2\\ 3 & 4\end{bmatrix} \in M_{2\times 2}(\bbR)$.

\begin{enumerate}

\item(\emph{1 point}) Let $t$ be a formal variable. Find the characteristic polynomial $\det(A-tI_2)$ and the eigenvalues of $A$.
\item[] $CP=(-1-t)(4-t)-(2)(3)=t^2-3t-10\Rightarrow(t-5)(t+2)$. Eigenvalues are 5 and $-2$.

\item (\emph{1 point}) Use your answer in the previous part to find the eigenvectors of $A$ of each eigenvalue.
\item[] $(A-5I)v=0$: $\begin{bmatrix}
    -6&2\\3&-1
\end{bmatrix}\Rightarrow\begin{cases}
    -6x_1+2x_2=0\\3x_1-1x_2=0
\end{cases}$, $x_2=3x_1\Rightarrow v_5=\text{span}\{(1,3)^T\}$.
\item[] $(A+2I)v=0$: $\begin{bmatrix}
    1&2\\3&6
\end{bmatrix}\Rightarrow\begin{cases}
    1x_1+2x_2=0\\3x_1+6x_2=0
\end{cases}$, $x_1=-2x_2\Rightarrow v_{-2}=\text{span}\{(-2,1)^T\}$.

\item (\emph{1 point}) Use your answer in the previous part to find an invertible matrix $C \in M_{2\times 2}(\bbR)$ and a diagonal matrix $B \in M_{2\times 2}(\bbR)$ such that $A=CBC^{-1}$
\item[] $C=\begin{bmatrix}
    1&-2\\3&1
\end{bmatrix}$, $B=\begin{bmatrix}
    5&0\\0&-2
\end{bmatrix}$. $\det(C)=1\cdot1-3\cdot-2=7\neq0$, so $A=CBC^-1$.

\item (\emph{2 points}) Use your answer to part $A$ to calculate $A^{51}$. (For clarity, you can leave any exponentials written as exponentials.)
\item[] $A^{51}=CB^{51}C^{-1}=\frac{1}{7}\begin{bmatrix}
    1&-2\\3&1
\end{bmatrix}\begin{bmatrix}
    5^{52}&0\\0&(-2)^{52}
\end{bmatrix}\begin{bmatrix}
    1&2\\-3&1
\end{bmatrix}$. Set $a=5^{51}$, $b=(-2)^{51}$.
\item[] $A^{51}=\frac{1}{7}\begin{bmatrix}
    a&-2b\\3a&b
\end{bmatrix}\begin{bmatrix}
    1&2\\-3&1
\end{bmatrix}=\frac{1}{7}\begin{bmatrix}
    a+6b&2a-2b\\3a-3b&6a+b
\end{bmatrix}$.

\item (\emph{2 points}) Find all ``square roots'' of $A$; i.e., find all (possibly complex) matrices $B$ such that $B^2=A$.
\item[] $B=C\begin{bmatrix}
    \pm\sqrt{5}&0\\0&\pm i\sqrt{2}
\end{bmatrix}C^{-1}=\frac{1}{7}\begin{bmatrix}
    1&-2\\3&1
\end{bmatrix}\begin{bmatrix}
    \pm\sqrt{5}&0\\0&\pm i\sqrt{2}
\end{bmatrix}\begin{bmatrix}
    1&2\\-3&1
\end{bmatrix}$. 

\end{enumerate}

\item

Let $L \colon \mathbb{C}^n \to \mathbb{C}^n$ be the linear transformation that cyclically permutes the coefficients of the vector: 
\[
L(a_1,a_2,\dots,a_n) = (a_2, \dots, a_{n}, a_1).
\]
\begin{enumerate}

\item (\emph{3 points}) Find the eigenvectors and eigenvalues of $L$. \textit{Hint:} Use the Laplace expansion to compute the characteristic polynomial. 
\item[] The characteristic polynomial of $L$ is $t^n-1$ (the permutation matrix of an $n$-cycle). The eigenvalues are the $n$-th roots: $\lambda_k=e^{2\pi ik/n}, k=0,1,\ldots,n-1$. For each $k$, $v^k=(1,\lambda_k,\lambda_k^2,\ldots,\lambda_k^{n-1})^T$. $Lv^k=(\lambda_k,\lambda_k^2,\ldots,\lambda_k^{n-1},\lambda_k^n)^T=\lambda_k(1,\lambda_k,\ldots,\lambda_k^{n-1})^T=\lambda_kv^k$. Thus, $v^k$ is an eigenvector with eigenvector $\lambda_k$. $n$ vectors are ind and form basis for $\bbC^n$.

\item (\emph{1 point}) Find the matrix of $L$ in the standard basis of $\mathbb{C}^n$.
\item[] $L=\begin{bmatrix}
    0&1&0&\ldots&0\\0&0&1&\ldots&0\\ \vdots&\vdots&\vdots&\ddots&1\\1&0&0&\ldots&0
\end{bmatrix}$.

\item (\emph{3 points}) Let $A$ be the $n \times n$ matrix that has $2$ on the diagonal, and 1 on both sides of the diagonal in such a way that they wrap around. For instance, if $n=5$, then $A$ would be the matrix
\[
\begin{bmatrix}
2 & 1 & 0 & 0 & 1 \\
1 & 2 & 1 & 0 & 0 \\
0 & 1 & 2 & 1 & 0 \\
0 & 0 & 1 & 2 & 1 \\
1 & 0 & 0 & 1 & 2 
\end{bmatrix}.
\]
Find the eigenvalues and eigenvectors of the matrix $A \in \mathrm{Mat}_{n \times n}(\mathbb{C})$ (in general, not just in the case $n = 5$). Simplify your answer as much as possible. \textit{Hint:} Use the previous parts instead of computing the characteristic polynomial. 
\item[] $A=2I+L+L^{-1}\Rightarrow Av^k=(2+\lambda_k+\lambda_k^{-1})v^k$, Simplify using $\lambda_k+\lambda_k^{-1}=2\cos(2\pi k/n)$
\item[] $\Rightarrow \lambda_k=2+2\cos(2\pi k/2)$, $k=0,1,\ldots n-1$. Eigenvectors are the same $v^k$ as above.

\end{enumerate}

\item (\emph{3 points}) Recall that if $V$ is a vector space, then $V^*$ is the dual space of $V$, see Problem 1 in Homework 6. Let $T\colon V \to W$ be a linear transformation. The \emph{dual linear transformation} $T^*\colon W^* \to V^*$ is defined as $$T^*(\phi)=\phi \circ T. $$
In other words, $T^*(\phi) \in V^*$ is the linear functional such that $(T^*(\phi))(v)=\phi(T(v))$ for any $v\in V$.

Let $\mathcal{A}=\{v_1,\ldots,v_n\} \subset V$ and  $\mathcal{B}=\{w_1,\ldots,w_m\} \subset W$ be bases and let $\mathcal{A}^\vee=\{\phi_1,\ldots,\phi_n\} \subset V^*$ and $\mathcal{B}^\vee=\{\psi_1,\ldots,\psi_m\} \subset W^*$ be the dual bases as in Problem 1a in Homework 6. Show that
$$[T^*]_{\mathcal{A}^\vee\mathcal{B}^\vee} =([T]_{\mathcal{B}\mathcal{A}})^T \in M_{n\times m}(\bbF)$$
is the transpose matrix of $T$.
\item[] By definition of the dual map and evaluation map from HW6(b), for each dual basis element $\psi_i\in\mathcal{B}^\vee$ and each basis vector $v_j\in\mathcal{A}$, $(T^*(\psi_i))(v_j)=\psi_i(T(v_j))$. Now suppose $T(v_j)=\sum_{r=1}^mt_{rj}w_r$. Then $(T^*(\psi_i))(v_j)=\psi_i(\sum_{r=1}^mt_{rj}w_r)=\sum_{r=1}^mt_{rj}\psi_i(w_r)$. Because $\psi_i(w_r)=\delta_{ir}$, only $r=i$ term remains, leaving $(T^*(\psi_i))(v_j)=t_{ij}$. This links entries of $[T]_{\mathcal{BA}}$ and $T^*(\psi_i)$.
\item[] Since $\{\phi_1,\ldots,\phi_n\}$ is a basis of $V^*$, there are scalars ($s$), with $T^*(\psi_i)=\sum_{k=1}^ns_{ki}\phi_k$. Using $\psi_k(v_j)=\delta_{kj}$ from HW6(a), $(T^*(\psi_i))(v_j)=\sum_{k=1}^ns_{ki} \phi_k(v_j)=\sum_{k=1}^ns_{ki}\delta_{kj}=s_{ji}$. However, from deduction above, we must have for every $i=1,\ldots,m$ and $j=1,\ldots,n$ that $s_{ji}=t_{ij}$. 
\item[] The matrix of $T$ w.r.t bases $\mathcal{B,A}$ is $[T]_{\mathcal{BA}}=(t_{ij})$. The matrix of $T^*$ w.r.t bases $\mathcal{B^\vee,A^\vee}$ is $[T^*]_{\mathcal{B^\vee A\vee}}=(s_{ji})$. From the equality above, this is the transpose relation. 
\item[] Thus $[T^*]_{\mathcal{A}^\vee\mathcal{B}^\vee} =([T]_{\mathcal{B}\mathcal{A}})^T$.
\end{enumerate}

\end{document}
